\scanpage{39}
\mspage{37}
% archive sheet number: 38
\deleted{de 8.1. Dans le \unsure{deuxième cas}, \inserted{cas d'un corps quelconque} \unsure{pour des généralités} \unsure{dans une} $V\otimes_{k_0}k'_0$, où $k'_0$ est une extension \unsure{radicielle|radicale} convenable de $k_0$, $k_0=\underset{\longrightarrow}{\lim}\,k_\alpha$, où $k_\alpha$ parcourt les sous-corps de $k_0$ qui sont des extensions de type fini \unsure{dans le} corps premier. Pour $\alpha$ assez grand, $S$ provient de $S_\alpha$ \inserted{lisse} sur $k_\alpha$, \inserted{et le diviseur $(t)$ de $S$ provient d'un diviseur \unsure{réduit} $D_\alpha$ sur $X_\alpha$.} Alors $V$ est limite des $\unsure{\mathscr O_{S_\alpha,x_\alpha}|\mathcal O_{S_\alpha,x_\alpha}}$, $x_\alpha$ étant \unsure{le point considéré dans $S_\alpha$} \inserted{le point générique de $D_\alpha$}. Or le \unsure{corps} $\mathscr O_{\unsure{S_\alpha|S_0},x_\alpha}$ est un anneau de valuation discrète à corps résiduel de type fini (au sens absolu).}
% slanted
\margin{\deleted{(ce qui ne change pas les groupes d'inertie), on est ramené au cas où $k(V)$ sur $k$ est \uline{séparable}. On peut donc supposer $S$ \uline{lisse} sur $k_0$, et \unsure{Écrivons}.}}

Dans \unsure{donc} les deux cas, l'anneau \deleted{\unsure{ainsi}} \inserted{(resp. $V\otimes_{k_0}k'_0$)} est \uline{lim} filtrante \inserted{(par morphismes d'inclusions locaux)} de \deleted{ses} sous-anneaux de valuation discrets, \inserted{$V_\alpha$ à corps résiduel de type fini. Donc,} \deleted{\unsure{avec corps rés. de type fini de $k$ considéré}.} le hensélisé $\unsure{V^h|V^{sh}}$ (resp. $(V\otimes_{k_0}k'_0)^h=V^h\otimes_{k_0}k'_0$) est limite inductive des $V_\alpha^h$, et on peut lui appliquer 8.1. De même pour le hensélisé strict. \deleted{\unsure{8}}

b) Cas $V$ de car. nulle, \inserted{i.e. algèbre sur $\unsure{\mathbb Q|Q}$.} $V$ est limite inductive de \inserted{ses} sous-$\mathbb Q$-algèbres locales $V_\alpha$, ess. de type fini sur $\mathbb Q$, \deleted{avec $V_\alpha$} d'idéaux maximaux $\unsure{\mathfrak m|m}\cap V_\alpha=\mathfrak m_\alpha$. On peut supposer $t\in V_\alpha$. Par Hironaka, on peut trouver un \deleted{\unsure{modèle normal}} schéma propre, \deleted{sur} intègre, \uuline{régulier} $S'_\alpha$ sur $S_\alpha=\operatorname{Spec}V_\alpha$, avec $S'_\alpha\to S_\alpha$ birationnel, et l'image inverse \inserted{dans $S'_\alpha$} \deleted{est} de $\operatorname{div}(t_{S_\alpha})$ un diviseur à croisements normaux. \deleted{\unsure{Prenant}} Comme $S'_\alpha-S_\alpha$ est propre, il existe un point $x'_\alpha$ de $S'_\alpha$ tel que l'anneau local $V'_\alpha\unsure{:|=}\mathscr O_{\unsure{S'_\alpha|\alpha'},x'_\alpha}$ soit dominé par $V$. \deleted{Donc} Or cet anneau local est régulier, et le diviseur de $t$ y est à croisements normaux.
% slanted
\margin{\deleted{\unsure{L'hypothèse}} donc $t=\prod x_i^{n_i}$, \unsure{les $x_i\in\mathfrak m'_\alpha$|les $x_i\in\mathfrak m_\alpha$}, $n_i\geq1$. D'où le fait qu'il n'y a qu'un seul $x_i$, et que \deleted{\unsure{de plus on a}} $\unsure{n_1|n_i}=1$.}
\deleted{Donc} les $V_\alpha$ ayant cette propriété supplémentaire
